\documentclass[11pt,aspectratio=169]{beamer}
% Stile comune delle slide del laboratorio (incluso da slides/it e slides/en)
\usetheme{default}
\usefonttheme{professionalfonts}
\setbeamertemplate{navigation symbols}{}
\setbeamertemplate{footline}[frame number]

\usepackage{amsmath,amssymb,booktabs}
\usepackage{eurosym}
\usepackage{tikz}
\usepackage{pgfplots}
\pgfplotsset{compat=1.16}
\usetikzlibrary{shapes.geometric,arrows.meta}

\definecolor{blunotte}{HTML}{16324A}
\definecolor{teal}{HTML}{0E7490}
\definecolor{tealchiaro}{HTML}{E3F2F5}
\definecolor{verde}{HTML}{1E8449}
\definecolor{arancio}{HTML}{CA6F1E}
\definecolor{rossomattone}{HTML}{C0392B}
\definecolor{grigiomedio}{HTML}{7F8C8D}
\definecolor{viola}{HTML}{8E44AD}
\definecolor{ocra}{HTML}{B7950B}

\setbeamercolor{structure}{fg=teal}
\setbeamercolor{frametitle}{fg=white, bg=blunotte}
\setbeamercolor{title}{fg=white, bg=blunotte}
\setbeamercolor{block title}{fg=white, bg=teal}
\setbeamercolor{block body}{bg=tealchiaro}
\setbeamercolor{alerted text}{fg=rossomattone}
\setbeamertemplate{itemize items}[circle]

\pgfplotsset{
  slide/.style={width=0.72\textwidth, height=5.4cm, grid=major,
    grid style={black!12}, tick label style={font=\scriptsize},
    label style={font=\small}, title style={font=\small\bfseries, color=blunotte},
    legend style={font=\scriptsize, draw=none, fill=none,
      at={(0.5,-0.28)}, anchor=north}, legend columns=3, legend cell align=left},
}
% i CSV delle figure vengono dalla dispensa: un'unica fonte di verità
\newcommand{\figdat}[1]{../../dispensa/figure/dat/#1.csv}
\newcommand{\R}{\mathbb{R}}
\newcommand{\E}{\mathbb{E}}
\newcommand{\vet}[1]{\boldsymbol{#1}}
\newcommand{\Prob}{\mathbb{P}}
\newcommand{\var}{\mathrm{VaR}}
\newcommand{\cvar}{\mathrm{CVaR}}
\newcommand{\T}{^{\mathsf T}}

\usepackage[italian]{babel}

\title{Laboratorio di Ricerca Operativa\\
  \large Modelli continui di ottimizzazione per l'Ingegneria Gestionale}
\author{Fabio Furini}
\institute{Sapienza Università di Roma}
\date{A.A. 2026--2027}

\AtBeginSection[]{
  \begin{frame}
    \vfill\centering
    \usebeamerfont{title}\usebeamercolor[fg]{structure}\insertsection
    \vfill
  \end{frame}}

\begin{document}

\begin{frame}
  \titlepage
\end{frame}

\begin{frame}{Indice del corso}
\small
\begin{columns}[t]
\column{0.5\textwidth}
\textbf{Strumenti}
\begin{itemize}\itemsep0pt
  \item Introduzione e filosofia del corso
  \item Richiami: LP, dualità, KKT, sensitività
  \item Il solver: costruire, risolvere, interpretare
\end{itemize}
\textbf{Modelli deterministici}
\begin{itemize}\itemsep0pt
  \item Produzione e scorte (LP/QP)
  \item Supply chain e CO$_2$ (LP/NLP)
  \item Portafoglio di Markowitz (QP)
  \item Pricing (NLP)
\end{itemize}
\column{0.5\textwidth}
\begin{itemize}\itemsep0pt
  \item Budget pubblicitario (NLP convesso)
  \item Localizzazione continua (NLP convesso)
  \item Ricarica di veicoli elettrici (LP/QP)
  \item Code e capacità (NLP convesso)
\end{itemize}
\textbf{Incertezza e machine learning}
\begin{itemize}\itemsep0pt
  \item Newsvendor (LP stocastico)
  \item VaR e CVaR (LP)
  \item Support Vector Machine (QP)
\end{itemize}
\textbf{Organizzazione del laboratorio}
\end{columns}
\end{frame}

% ==================================================================
\section{Introduzione}

\begin{frame}{Di che cosa si occupa questo laboratorio}
\begin{itemize}
  \item Trasformare \alert{decisioni aziendali} in modelli di ottimizzazione:
        quanto produrre, dove localizzare, che prezzo fissare, quanto rischio accettare.
  \item Implementare i modelli in \textbf{Python + Gurobi}, risolverli e --- soprattutto ---
        \alert{interrogarli}: quanto vale un'ora di capacità in più? la soluzione regge se
        i dati cambiano del 5\%?
  \item Solo \textbf{variabili continue}: valgono dualità, prezzi ombra e condizioni KKT ---
        gli strumenti che trasformano numeri in spiegazioni economiche.
\end{itemize}
\begin{block}{La domanda finale di ogni esercitazione}
Non solo \emph{``qual è l'ottimo?''}, ma \emph{``quale decisione suggeriamo e quanto è robusta?''}
\end{block}
\end{frame}

\begin{frame}{Notazione e sigle}
\begin{itemize}
  \item \textbf{LP} programmazione lineare; \textbf{QP} quadratica; \textbf{NLP} non lineare.
        \emph{Convesso} $=$ ogni minimo locale è globale.
  \item Scalari e indici minuscoli ($x_{it}$, $\lambda$); \textbf{vettori} minuscoli in
        grassetto ($\vet x$); \textbf{matrici} maiuscole in grassetto ($\vet Q$);
        \textbf{insiemi} maiuscoli ($I$, $T$, $S$).
  \item Eccezioni (per aderenza alla letteratura): variabili aleatorie $D$, $L$;
        funzione di ripartizione $F$; parametro $C$ della SVM.
  \item In ogni modello: prima i \alert{dati} (con il loro tipo), poi le
        \alert{variabili}, poi obiettivo e vincoli spiegati riga per riga.
\end{itemize}
\end{frame}

% ==================================================================
\section{Richiami: LP, dualità, KKT}

\begin{frame}{La coppia primale-duale}
\begin{block}{LP in forma generale (per componenti)}
\vspace{-2ex}
\begin{align*}
\text{(P)}\;\; & \min \sum_{j=1}^{n} c_j x_j
&& \text{s.t.}\; \sum_{j=1}^{n} a_{kj} x_j \ge b_k \;\forall k, \quad x_j \ge 0\\
\text{(D)}\;\; & \max \sum_{k=1}^{m} b_k y_k
&& \text{s.t.}\; \sum_{k=1}^{m} a_{kj} y_k \le c_j \;\forall j, \quad y_k \ge 0
\end{align*}
\end{block}
\begin{itemize}
  \item \textbf{Dualità forte}: all'ottimo (P) e (D) hanno lo stesso valore.
  \item $y_k$ = \alert{prezzo ombra} della risorsa $k$: variazione dell'ottimo per
        una unità in più di $b_k$.
\end{itemize}
\end{frame}

\begin{frame}{Dualità su un esempio $2\times 2$, tutto a mano}
$\max 30x_A + 50x_B$ \; s.t. \; $x_A + 3x_B \le 90$ (ore), \; $2x_A + x_B \le 80$ (kg).
\begin{enumerate}
  \item Entrambi i vincoli attivi al vertice ottimo: $x_B = 20$, $x_A = 30$,
        $z^* = 1900$~\euro.
  \item Duali: $y_1 + 2y_2 = 30$, $3y_1 + y_2 = 50$ $\Rightarrow$
        $y_1 = 14$~\euro/ora, $y_2 = 8$~\euro/kg.
  \item Verifica: $90 \cdot 14 + 80 \cdot 8 = 1900 = z^*$. \checkmark
\end{enumerate}
\begin{block}{Analisi di sensitività}
Ore extra a 12~\euro/ora? Conviene ($14 > 12$). Con 100 ore: ricavo $2040 = 1900 +
10 \cdot 14$ --- il prezzo ombra è esatto finché le ore restano in $[40, 240]$
(\texttt{SARHSLow/Up}). Un vincolo non attivo ha sempre prezzo ombra nullo.
\end{block}
\end{frame}

\begin{frame}{Condizioni KKT}
Per $\min f(\vet x)$ s.t. $g_i(\vet x) \le 0$, $h_j(\vet x) = 0$, in un ottimo
regolare esistono $\lambda_i \ge 0$, $\nu_j$:
\[
\nabla f + \textstyle\sum_i \lambda_i \nabla g_i + \sum_j \nu_j \nabla h_j = \vet 0,
\qquad \lambda_i\, g_i(\vet x^*) = 0 \;\forall i \;\;\text{(complementarietà)}
\]
Se il problema è convesso, sono anche \alert{sufficienti}.
\begin{block}{Esempio svolto}
$\min x^2 + y^2$ s.t. $x + y \ge 4$: stazionarietà $x = y = \lambda/2$; vincolo
attivo $\Rightarrow x = y = 2$, $\lambda = 4$. Con termine noto $4 + \varepsilon$:
$f^* = 8 + 4\varepsilon + \varepsilon^2/2$ --- il moltiplicatore è la derivata
dell'ottimo.
\end{block}
\begin{itemize}
  \item Quando il solver non dà i duali: stimarli \alert{per perturbazione}.
\end{itemize}
\end{frame}

\begin{frame}{Il protocollo di analisi di sensitività (in ogni esercitazione)}
\begin{enumerate}
  \item \textbf{Scenario base}: risolvere, verificare, vincoli attivi.
  \item \textbf{One-at-a-time}: un parametro chiave su una griglia.
  \item \textbf{Prezzi ombra}: duale vs ri-ottimizzazione perturbata.
  \item \textbf{Scenari}: pessimistico, centrale, ottimistico.
  \item \textbf{Trade-off}: una frontiera (costo-servizio, rischio-rendimento, costo-CO$_2$).
  \item \textbf{Stabilità}: dati $\pm 5\%$ $\to$ la raccomandazione regge?
\end{enumerate}
\end{frame}

% ==================================================================
\section{Il solver: costruire, risolvere, interpretare}

\begin{frame}[fragile]{Costruire un modello in cinque passi}
\begin{semiverbatim}\small
m = gp.Model("produzione")                   \textcolor{grigiomedio}{# 1. contenitore}
x = m.addVars(prodotti, mesi, name="x")      \textcolor{grigiomedio}{# 2. variabili (>= 0 di default)}
m.addConstrs((x.sum("*", t) <= b[t]          \textcolor{grigiomedio}{# 3. vincoli}
              for t in mesi), name="cap")
m.setObjective(..., GRB.MINIMIZE)            \textcolor{grigiomedio}{# 4. obiettivo}
m.write("modello.lp"); m.optimize()          \textcolor{grigiomedio}{# 5. debug + soluzione}
\end{semiverbatim}
\begin{itemize}
  \item Variabili \alert{libere} (SVM: $b$; CVaR: $\eta$):
        \texttt{m.addVar(lb=-GRB.INFINITY)} --- dimenticarlo è l'errore più comune!
  \item Somme: \texttt{gp.quicksum(...)} o \texttt{x.sum(i, "*")}.
\end{itemize}
\end{frame}

\begin{frame}[fragile]{Leggere la soluzione}
\begin{semiverbatim}\small
if m.Status == GRB.OPTIMAL:      \textcolor{grigiomedio}{# SEMPRE prima di leggere numeri}
    print(m.ObjVal, x[i,t].X)
elif m.Status == GRB.INFEASIBLE:
    m.computeIIS(); m.write("conflitto.ilp")
\end{semiverbatim}
\begin{center}\small
\begin{tabular}{ll}
\toprule
\texttt{Pi} & prezzo ombra del vincolo (LP) \\
\texttt{Slack} & scarto: $0$ $\Rightarrow$ vincolo attivo \\
\texttt{RC} & costo ridotto di una variabile a zero \\
\texttt{SARHSLow/Up} & intervallo di validità del prezzo ombra \\
\bottomrule
\end{tabular}
\end{center}
\begin{itemize}
  \item Problemi di minimo, vincolo $\le$: \texttt{Pi} $\le 0$ (convenzione Gurobi).
  \item RHS con costanti a sinistra: Gurobi le sposta nel termine noto
        (\alert{trappola} nell'analisi di sensitività --- v.\ ricarica EV).
\end{itemize}
\end{frame}

\begin{frame}{Checklist di interpretazione}
\begin{enumerate}
  \item Stato \texttt{OPTIMAL}? Se no: diagnosi (IIS, file \texttt{.lp}).
  \item Ordine di grandezza dell'ottimo sensato?
  \item Vincoli attivi $=$ quelli attesi?
  \item Prezzi ombra: quale risorsa potenziare per prima, e fino a che prezzo?
  \item Stabilità: dati $\pm 5\%$.
  \item Raccomandazione manageriale in tre righe.
\end{enumerate}
\begin{block}{Per gli NLP generali}
\texttt{scipy.optimize.minimize} (SLSQP, Nelder-Mead): stessa logica, ma ottimo
\alert{locale} salvo convessità dimostrata.
\end{block}
\end{frame}

% ==================================================================
\section{Produzione e scorte multiperiodali (LP/QP)}

\begin{frame}{Produzione e scorte --- il problema}
\textbf{A parole.} Decidere quanto produrre e quanto tenere a scorta, mese per mese.
Obiettivo: minimo costo di produzione $+$ giacenza. Vincoli: bilancio di magazzino
(domanda servita per intero) e ore macchina disponibili.
\begin{block}{Dati e variabili}
Dati: domanda $d_{it}$, costi $c_{it}, h_{it}$, ore per unità $a_i$, ore disponibili
$b_t$, scorta iniziale $s_{i0}$ (tutti reali $\ge 0$).\\
Variabili continue: $x_{it} \ge 0$ produzione, $s_{it} \ge 0$ scorta di fine mese.
\end{block}
\begin{block}{Modello (LP)}
\vspace{-2ex}
\begin{align*}
\min\;& \sum_{i \in I}\sum_{t \in T} ( c_{it} x_{it} + h_{it} s_{it} ) \\
\text{s.t.}\;\; & s_{i,t-1} + x_{it} = d_{it} + s_{it} \;\; \forall i \in I,\ t \in T
\qquad \textstyle\sum_{i \in I} a_i x_{it} \le b_t \;\; \forall t \in T
\end{align*}
\end{block}
\end{frame}

\begin{frame}{Produzione e scorte --- capire i vincoli}
\begin{itemize}
  \item \textbf{Obiettivo}: il modello scambia liberamente produzione anticipata
        (più giacenza) con produzione al limite di capacità.
  \item \textbf{Bilancio}: contabilità di magazzino; è il vincolo che \alert{lega i
        mesi}: senza, il problema si spezza in $|T|$ problemi indipendenti.
  \item \textbf{Capacità}: la risorsa scarsa condivisa; il suo duale dice quanto
        vale un'ora in più.
\end{itemize}
\begin{block}{Esempio a mano (1 prodotto, 2 mesi)}
$d = (60, 100)$, capacità $80$/mese, $c = 10$, $h = 1$: unico piano ammissibile
$x = (80, 80)$, $s_1 = 20$, costo $1620$~\euro. $+1$ ora nel mese 2 $\Rightarrow$
costo $-1$~\euro{} (un mese di giacenza in meno): duale mese 2 $= -1$, mese 1 $= 0$.
\end{block}
\end{frame}

\begin{frame}{Produzione e scorte --- caso di studio}
3 prodotti $\times$ 6 mesi, picco di domanda nei mesi 4--5, 420--460 ore/mese.
\begin{columns}
\column{0.62\textwidth}
\begin{tikzpicture}
\begin{axis}[slide, width=\textwidth, ybar, bar width=4pt, xtick={1,...,6},
  xlabel={mese}, ylabel={unità}, area legend, enlarge x limits=0.1, ymin=0]
\addplot [fill=teal!85, draw=none] table [col sep=comma, x=mese, y=x1] {\figdat{cap04_piano}};
\addplot [fill=rossomattone!80, draw=none] table [col sep=comma, x=mese, y=x2] {\figdat{cap04_piano}};
\addplot [fill=arancio!85, draw=none] table [col sep=comma, x=mese, y=x3] {\figdat{cap04_piano}};
\addplot [sharp plot, grigiomedio, thick, dashed, mark=o]
  table [col sep=comma, x=mese, y=domtot] {\figdat{cap04_piano}};
\legend{prod.\ 1, prod.\ 2, prod.\ 3, domanda tot.}
\end{axis}
\end{tikzpicture}
\column{0.38\textwidth}\small
\begin{itemize}
  \item Costo ottimo: \textbf{32.889~\euro}
  \item Il prodotto 3 (2,1 ore/unità) viene prodotto \alert{in anticipo}
        (\emph{pre-build}) e messo a scorta (fino a 107 unità) per liberare i mesi
        di picco.
  \item Scorte a zero nel mese finale: il modello non spreca.
\end{itemize}
\end{columns}
\end{frame}

\begin{frame}{Produzione e scorte --- i prezzi ombra raccontano una storia}
\begin{center}\small
\begin{tabular}{lcccccc}
\toprule
mese & 1 & 2 & 3 & 4 & 5 & 6\\
\midrule
uso ore & 330/420 & 420/420 & 460/460 & 460/460 & 460/460 & 420/420 \\
prezzo ombra (\euro/h) & 0 & $-0{,}76$ & $-1{,}52$ & $-2{,}29$ & $-3{,}05$ & $-3{,}81$ \\
\bottomrule
\end{tabular}
\end{center}
\begin{itemize}
  \item I duali crescono di $0{,}762 = h_3/a_3$ al mese: un'ora in più nel mese $t$
        fa risparmiare \alert{un mese di giacenza} del prodotto 3.
  \item Verifica per perturbazione: $+1$ ora nel mese 6 $\Rightarrow$ costo
        $-3{,}810$~\euro{} (coincide col duale).
  \item Quando solver e ragionamento economico danno lo stesso numero, il modello è
        probabilmente giusto.
\end{itemize}
\end{frame}

\begin{frame}{Produzione e scorte --- smoothing e valore della capacità}
\begin{columns}
\column{0.55\textwidth}
\begin{tikzpicture}
\begin{axis}[slide, width=\textwidth, xlabel={capacità (\% dell'attuale)},
  ylabel={costo ottimo (\euro)}]
\addplot+ [mark=*, mark size=1.3pt] table [col sep=comma, x=percento, y=costo]
  {\figdat{cap04_capacita}};
\end{axis}
\end{tikzpicture}
\column{0.45\textwidth}\small
\begin{itemize}
  \item Curva \textbf{convessa e decrescente}: le prime ore extra valgono molto.
  \item Sotto il 97\%: \alert{inammissibile} --- la domanda non è più servibile.
  \item Variante QP con $\gamma \sum_{t>1} (X_t - X_{t-1})^2$: profilo quasi piatto
        (variazione max da 81 a 1 unità) per $+114$~\euro{} ($+0{,}35\%$).
\end{itemize}
\end{columns}
\begin{block}{Take-away}
Ore straordinarie nei mesi 5--6 convengono fino a 3--3,8~\euro/ora sotto il prezzo
interno; il piano livellato è quasi gratis.
\end{block}
\end{frame}

% ==================================================================
\section{Supply chain con congestione e CO$_2$ (LP/NLP)}

\begin{frame}{Supply chain --- il problema}
\textbf{A parole.} Decidere quante unità far viaggiare su ogni tratta della rete.
Obiettivo: minimo costo di trasporto. Vincoli: conservazione del flusso in ogni nodo
e capacità delle tratte.
\begin{block}{Modello (flusso a costo minimo)}
\vspace{-2ex}
\[
\min \sum_{(i,j) \in A} c_{ij} x_{ij}
\;\;\text{s.t.}\;
\sum_{j:(i,j) \in A} x_{ij} - \sum_{j:(j,i) \in A} x_{ji} = b_i \;\forall i \in N,
\;\; 0 \le x_{ij} \le u_{ij}
\]
\end{block}
Varianti sull'obiettivo: congestione quadratica
$\sum (c_{ij} x_{ij} + \alpha c_{ij} x_{ij}^2/u_{ij})$ (convessa); prezzo interno
CO$_2$: $\sum (c_{ij} + \tau e_{ij}) x_{ij}$; minimax dell'utilizzo.
\begin{itemize}
  \item Esempio a mano (2 rotte): l'LP satura la rotta economica; la congestione
        eguaglia i \alert{costi marginali} e lascia margine di sicurezza.
\end{itemize}
\end{frame}

\begin{frame}{Supply chain --- caso di studio}
\begin{columns}
\column{0.52\textwidth}
\centering
\scalebox{0.8}{% Rete della soluzione: LP a costo minimo (generato da lab05_supplychain.py)
\begin{tikzpicture}[>=stealth,
    nodo/.style={circle, draw=none, text=white, font=\bfseries\small,
                 minimum size=8mm, inner sep=0pt}]
  \draw[rossomattone, line width=2.00pt] (0.00,1.15) -- (3.40,0.92)
    node[midway, above, sloped, font=\tiny, text=black!60] {220};
  \draw[teal, line width=0.47pt] (0.00,1.15) -- (3.40,-0.92)
    node[midway, above, sloped, font=\tiny, text=black!60] {10};
  \draw[black!25, densely dotted, line width=0.4pt] (0.00,-1.15) -- (3.40,0.92);
  \draw[rossomattone, line width=2.00pt] (0.00,-1.15) -- (3.40,-0.92)
    node[midway, above, sloped, font=\tiny, text=black!60] {220};
  \draw[teal, line width=1.27pt] (3.40,0.92) -- (6.80,1.72)
    node[midway, above, sloped, font=\tiny, text=black!60] {120};
  \draw[teal, line width=1.05pt] (3.40,0.92) -- (6.80,0.57)
    node[midway, above, sloped, font=\tiny, text=black!60] {90};
  \draw[teal, line width=0.47pt] (3.40,0.92) -- (6.80,-0.57)
    node[midway, above, sloped, font=\tiny, text=black!60] {10};
  \draw[black!25, densely dotted, line width=0.4pt] (3.40,0.92) -- (6.80,-1.72);
  \draw[black!25, densely dotted, line width=0.4pt] (3.40,-0.92) -- (6.80,1.72);
  \draw[black!25, densely dotted, line width=0.4pt] (3.40,-0.92) -- (6.80,0.57);
  \draw[rossomattone, line width=1.35pt] (3.40,-0.92) -- (6.80,-0.57)
    node[midway, above, sloped, font=\tiny, text=black!60] {130};
  \draw[teal, line width=1.13pt] (3.40,-0.92) -- (6.80,-1.72)
    node[midway, above, sloped, font=\tiny, text=black!60] {100};
  \node[nodo, fill=verde] at (0.00,1.15) {S1};
  \node[nodo, fill=verde] at (0.00,-1.15) {S2};
  \node[nodo, fill=arancio] at (3.40,0.92) {H1};
  \node[nodo, fill=arancio] at (3.40,-0.92) {H2};
  \node[nodo, fill=teal] at (6.80,1.72) {M1};
  \node[nodo, fill=teal] at (6.80,0.57) {M2};
  \node[nodo, fill=teal] at (6.80,-0.57) {M3};
  \node[nodo, fill=teal] at (6.80,-1.72) {M4};
  \node[font=\small\bfseries, text=blunotte] at (3.40,2.30) {LP a costo minimo};
\end{tikzpicture}}\\[0.5ex]
\scalebox{0.8}{% Rete della soluzione: Congestione quadratica (generato da lab05_supplychain.py)
\begin{tikzpicture}[>=stealth,
    nodo/.style={circle, draw=none, text=white, font=\bfseries\small,
                 minimum size=8mm, inner sep=0pt}]
  \draw[teal, line width=1.65pt] (0.00,1.15) -- (3.40,0.92)
    node[midway, above, sloped, font=\tiny, text=black!60] {172};
  \draw[teal, line width=0.73pt] (0.00,1.15) -- (3.40,-0.92)
    node[midway, above, sloped, font=\tiny, text=black!60] {45};
  \draw[teal, line width=0.65pt] (0.00,-1.15) -- (3.40,0.92)
    node[midway, above, sloped, font=\tiny, text=black!60] {35};
  \draw[teal, line width=1.83pt] (0.00,-1.15) -- (3.40,-0.92)
    node[midway, above, sloped, font=\tiny, text=black!60] {197};
  \draw[teal, line width=1.16pt] (3.40,0.92) -- (6.80,1.72)
    node[midway, above, sloped, font=\tiny, text=black!60] {105};
  \draw[teal, line width=0.68pt] (3.40,0.92) -- (6.80,0.57)
    node[midway, above, sloped, font=\tiny, text=black!60] {39};
  \draw[teal, line width=0.69pt] (3.40,0.92) -- (6.80,-0.57)
    node[midway, above, sloped, font=\tiny, text=black!60] {40};
  \draw[teal, line width=0.57pt] (3.40,0.92) -- (6.80,-1.72)
    node[midway, above, sloped, font=\tiny, text=black!60] {23};
  \draw[teal, line width=0.51pt] (3.40,-0.92) -- (6.80,1.72)
    node[midway, above, sloped, font=\tiny, text=black!60] {15};
  \draw[teal, line width=0.77pt] (3.40,-0.92) -- (6.80,0.57)
    node[midway, above, sloped, font=\tiny, text=black!60] {51};
  \draw[teal, line width=1.12pt] (3.40,-0.92) -- (6.80,-0.57)
    node[midway, above, sloped, font=\tiny, text=black!60] {100};
  \draw[teal, line width=0.96pt] (3.40,-0.92) -- (6.80,-1.72)
    node[midway, above, sloped, font=\tiny, text=black!60] {77};
  \node[nodo, fill=verde] at (0.00,1.15) {S1};
  \node[nodo, fill=verde] at (0.00,-1.15) {S2};
  \node[nodo, fill=arancio] at (3.40,0.92) {H1};
  \node[nodo, fill=arancio] at (3.40,-0.92) {H2};
  \node[nodo, fill=teal] at (6.80,1.72) {M1};
  \node[nodo, fill=teal] at (6.80,0.57) {M2};
  \node[nodo, fill=teal] at (6.80,-0.57) {M3};
  \node[nodo, fill=teal] at (6.80,-1.72) {M4};
  \node[font=\small\bfseries, text=blunotte] at (3.40,2.30) {Congestione quadratica};
\end{tikzpicture}}
\column{0.48\textwidth}\small
2 stabilimenti $\to$ 2 hub $\to$ 4 mercati; tratte economiche ``su strada''
(inquinanti), care ``su ferro'' (pulite).
\begin{itemize}
  \item LP: costo 3.385~\euro, \alert{3 archi saturi} (in rosso), emissioni 2.730 kg.
  \item Congestione: costo 3.703~\euro, nessun arco oltre il 90\%.
  \item Prezzi ombra della domanda $=$ costi marginali di servizio dei mercati:
        M1 8,00 --- M4 10,50 \euro/unità.
\end{itemize}
\end{columns}
\end{frame}

\begin{frame}{Supply chain --- il prezzo della CO$_2$}
\begin{columns}
\column{0.55\textwidth}
\begin{tikzpicture}
\begin{axis}[slide, width=\textwidth, xlabel={emissioni (kg CO$_2$)},
  ylabel={costo di trasporto (\euro)}]
\addplot+ [mark=*, mark size=1.5pt] table [col sep=comma, x=emissioni, y=costo]
  {\figdat{cap05_frontiera_co2}};
\end{axis}
\end{tikzpicture}
\column{0.45\textwidth}\small
\begin{itemize}
  \item $\tau \le 1$: assetto ``strada'' (3.385~\euro, 2.730 kg).
  \item $\tau \ge 1{,}5$: assetto ``ferro'' (4.705~\euro, 1.661 kg).
  \item Soglia esatta (bisezione): $\tau^* = 1{,}25$~\euro/kg.
  \item Le soluzioni LP stanno nei \alert{vertici}: la frontiera procede a salti;
        sotto la soglia il prezzo interno della CO$_2$ non cambia nulla.
\end{itemize}
\end{columns}
\begin{block}{Attenzione}
Il minimax puro ($\min z$) è degenere: qualunque instradamento con utilizzo
$\le z^*$ è ottimo, anche costosissimo --- combinare sempre con il costo.
\end{block}
\end{frame}

% ==================================================================
\section{Portafoglio di Markowitz (QP)}

\begin{frame}{Markowitz --- il problema}
\textbf{A parole.} Ripartire il capitale tra $n$ titoli. Obiettivo: minima varianza.
Vincoli: quote a somma 1, rendimento atteso minimo, limiti per titolo.
\begin{block}{Modello (QP convesso: $\vet Q \succeq 0$ sempre)}
\vspace{-2ex}
\[
\min \sum_{i \in I}\sum_{j \in I} q_{ij} x_i x_j \quad
\text{s.t.}\;\; \sum_{i \in I} \mu_i x_i \ge \bar r,\quad
\sum_{i \in I} x_i = 1,\quad \ell_i \le x_i \le u_i
\]
\end{block}
\begin{block}{Esempio a mano (2 titoli non correlati)}
$\sigma_1 = 20\%$, $\sigma_2 = 30\%$: $x_1 = 9/13 = 69{,}2\%$, volatilità di
portafoglio $16{,}6\%$ --- \alert{meno di entrambi i titoli}: la diversificazione
nella sua forma più pura.
\end{block}
\end{frame}

\begin{frame}{Markowitz --- la frontiera efficiente}
\begin{columns}
\column{0.62\textwidth}
\begin{tikzpicture}
\begin{axis}[slide, width=\textwidth, height=6.2cm,
  xlabel={volatilità annua (\%)}, ylabel={rendimento atteso (\%)}, xmin=4]
\addplot [teal, very thick] table [col sep=comma, x=vol, y=rend] {\figdat{cap06_front_libera}};
\addplot [arancio, thick, dashed] table [col sep=comma, x=vol, y=rend] {\figdat{cap06_front_limiti}};
\addplot [only marks, mark=*, mark size=1.6pt, grigiomedio]
  table [col sep=comma, x=vol, y=mu] {\figdat{cap06_titoli}};
\legend{senza limiti, $u_i = 30\%$, titoli singoli}
\end{axis}
\end{tikzpicture}
\column{0.38\textwidth}\small
8 ETF, $\mu$ e $\vet Q$ \alert{stimati} da 60 mesi.
\begin{itemize}
  \item Minima varianza: rendimento 11\%, volatilità \textbf{6\%} $<$ 7,5\% del
        miglior titolo.
  \item Equipesato $1/n$: \alert{dominato} (10,7\%, 10,1\%).
  \item Tetto $u_i = 30\%$: taglia la parte alta --- il costo dei mandati
        \emph{si vede}.
\end{itemize}
\end{columns}
\end{frame}

\begin{frame}{Markowitz --- sensitività e fragilità}
\begin{itemize}
  \item Vincolo di rendimento \alert{inattivo} fino a $\bar r = 11{,}06\%$ (la
        minima varianza rende già tanto): chiedere ``almeno l'8\%'' non costa nulla.
  \item Oltre: ogni punto di rendimento si paga in varianza (moltiplicatore
        $\approx 0{,}022$ a $\bar r = 12\%$).
  \item Composizione lungo la frontiera: dal mix difensivo (UTL, CON, SAN) al
        portafoglio concentrato sui titoli aggressivi (MAT, ENE).
\end{itemize}
\begin{block}{La lezione pratica più importante}
Le stime dei rendimenti sono \alert{molto più instabili} di quelle delle covarianze
(errore standard $\approx 2{,}6\%$ annuo su 60 mesi): i portafogli ottimizzati
inseguono gli errori di stima. Rimedi: vincoli $u_i$, shrinkage, minima varianza.
\end{block}
\end{frame}

% ==================================================================
\section{Pricing e revenue management (NLP)}

\begin{frame}{Pricing --- il problema}
\textbf{A parole.} Decidere prezzo $p$ e quantità $q$. Obiettivo: massimo profitto
$(p - c)q$. Vincoli: $q \le d(p)$ e $q \le k$ (capacità).
\begin{block}{Modello e domande}
\vspace{-1ex}
\[
\max\, p\,q - c\,q \;\;\text{s.t.}\;\; q \le d(p),\; q \le k
\]
lineare $d(p) = a - bp$; elasticità costante $d(p) = \theta p^{-\varepsilon}$;
logistica $d(p) = m/(1 + e^{-\alpha + \beta p})$.
\end{block}
\begin{itemize}
  \item La non convessità sta nel prodotto $p \cdot q$ (bilineare); con domanda
        lineare il profitto ridotto è una parabola concava; Gurobi risolve comunque
        al globale (\texttt{NonConvex=2}).
\end{itemize}
\end{frame}

\begin{frame}{Pricing --- esempio a mano e caso di studio}
Concerto: $d(p) = 1200 - 5p$, $c = 20$~\euro, $k = 400$ posti.
\begin{enumerate}
  \item Senza capacità: $p^\circ = (a/b + c)/2 = 130$~\euro, $q^\circ = 550$.
  \item La capacità morde: $p^* = (a - k)/b = \textbf{160}$~\euro,
        profitto 56.000~\euro.
  \item Valore di un posto in più: $(p^* - c) + k\,dp^*/dk = 140 - 80 =
        \textbf{60}$~\euro{} --- \alert{non} il margine pieno: per riempire il posto
        si abbassa il prezzo a tutti. (Gurobi, per perturbazione: 59,80.)
\end{enumerate}
\begin{block}{L'ottimo nel punto di spigolo}
Con elasticità costante l'ottimo non vincolato sarebbe $36{,}67$~\euro, ma lì la
domanda esplode: l'ottimo vero è nello spigolo $d(p) = k$, cioè
$p^* = (\theta/k)^{1/\varepsilon} = 79{,}11$~\euro{} --- nessuna derivata si annulla.
\end{block}
\end{frame}

\begin{frame}{Pricing --- profitto e capienza}
\begin{columns}
\column{0.52\textwidth}
\begin{tikzpicture}
\begin{axis}[slide, width=\textwidth, xlabel={prezzo (\euro)}, ylabel={profitto (\euro)}]
\addplot [grigiomedio, dashed, thick] table [col sep=comma, x=p, y=senza] {\figdat{cap07_profitto}};
\addplot [teal, very thick] table [col sep=comma, x=p, y=con] {\figdat{cap07_profitto}};
\addplot [rossomattone, dash dot] coordinates {(160, 0) (160, 60500)};
\legend{senza capienza, con $k = 400$, $p^* = 160$}
\end{axis}
\end{tikzpicture}
\column{0.48\textwidth}\small
\begin{itemize}
  \item Il vincolo di capienza taglia la parabola e sposta l'ottimo da 130 a
        160~\euro.
  \item La curva del valore della capienza si appiattisce a $q^\circ = 550$: oltre,
        i posti non valgono nulla.
  \item Multiprodotto con sostituzione (platea/galleria): prezzi da decidere
        \alert{congiuntamente} --- profitto 85.149~\euro; conviene ampliare la
        galleria ($+3.838$~\euro/50 posti), non la platea ($+3.287$).
\end{itemize}
\end{columns}
\end{frame}

% ==================================================================
\section{Budget pubblicitario (NLP convesso)}

\begin{frame}{Budget --- il problema e la condizione KKT}
\textbf{A parole.} Ripartire il budget $b$ tra canali con risposta concava
$r_i(x_i) = a_i \log(1 + k_i x_i)$; vincoli: budget e tetti $u_i$.
\[
\max \sum_{i \in I} r_i(x_i) \quad \text{s.t.}\;\; \sum_{i \in I} x_i \le b,
\quad 0 \le x_i \le u_i
\]
\begin{block}{KKT $=$ economia}
All'ottimo $r_i'(x_i^*) = \lambda$ per ogni canale interno:
\alert{l'ultimo euro rende lo stesso in tutti i canali attivi}. Canali al tetto:
marginale $> \lambda$; canali a zero: marginale iniziale $< \lambda$.
\end{block}
Esempio a mano ($b = 50$): $x = (35{,}6;\, 14{,}4)$, $\lambda = 2{,}46$.
\end{frame}

\begin{frame}{Budget --- verifica numerica e curva del valore}
\begin{center}\small
\begin{tabular}{lrrrr}
\toprule
 & social & search & TV & influencer \\
\midrule
spesa (migliaia \euro) & 24,3 & 34,8 & 22,9 & 18,0 \\
ritorno marginale & 8,2693 & 8,2693 & 8,2693 & 8,2693 \\
\bottomrule
\end{tabular}
\end{center}
\begin{itemize}
  \item Verifica: $+1000$~\euro{} di budget $\Rightarrow$ risposta $+8{,}244 \approx \lambda$.
  \item La TV riceve meno dei social nonostante il tetto più alto: conta la
        \alert{velocità di saturazione} $k_i$, non la dimensione del canale.
  \item $\lambda$ scende da 19 ($b = 20$) a 0 ($b = 300$, tutti ai tetti): la curva
        di $\lambda$ è l'argomento per negoziare il budget con il CFO.
\end{itemize}
\end{frame}

% ==================================================================
\section{Localizzazione continua (NLP convesso)}

\begin{frame}{Localizzazione --- tre obiettivi, tre risposte}
\textbf{A parole.} Decidere le coordinate $(x, y)$ della struttura; l'obiettivo
--- media, massimo o quadratica --- è la vera scelta manageriale.
\[
\text{Weber: } \min \textstyle\sum_{i \in I} w_i \sqrt{(x - a_i)^2 + (y - b_i)^2}
\qquad
\text{Minimax: } \min z, \; \text{dist}_i \le z \;\forall i \in I
\]
\begin{block}{Esempio a mano in 1D (pesi 1, 1, 3 nei km 0, 4, 10)}
Weber $=$ mediana pesata $= 10$; baricentro $= 6{,}8$; minimax $= 5$.
Tre obiettivi ragionevoli, tre risposte diverse.
\end{block}
\begin{itemize}
  \item Quadratico $\Rightarrow$ baricentro pesato (forma chiusa); Weber: non
        differenziabile nei punti di domanda; minimax: centro del cerchio minimo,
        \alert{ignora i pesi}.
\end{itemize}
\end{frame}

\begin{frame}{Localizzazione --- caso di studio ed equità}
\begin{columns}
\column{0.5\textwidth}
\centering
\scalebox{0.62}{% Mappa dei quartieri e localizzazioni ottime (generato da lab09_localizzazione.py)
\begin{tikzpicture}[scale=0.82, >=stealth]
  \draw[black!20, very thin] (0,0) grid[step=1] (10.5,10);
  \draw[->, black!50] (0,0) -- (10.8,0) node[below left, font=\scriptsize] {km est};
  \draw[->, black!50] (0,0) -- (0,10.3) node[below left, rotate=90, font=\scriptsize] {km nord};
  \fill[teal, opacity=0.45] (1.00,8.50) circle (0.38);
  \node[font=\tiny, text=black!55, anchor=west] at (1.38,8.50) {Q1};
  \fill[teal, opacity=0.45] (2.50,6.00) circle (0.47);
  \node[font=\tiny, text=black!55, anchor=west] at (2.97,6.00) {Q2};
  \fill[teal, opacity=0.45] (4.00,9.00) circle (0.33);
  \node[font=\tiny, text=black!55, anchor=west] at (4.33,9.00) {Q3};
  \fill[teal, opacity=0.45] (5.50,7.50) circle (0.52);
  \node[font=\tiny, text=black!55, anchor=west] at (6.02,7.50) {Q4};
  \fill[teal, opacity=0.45] (7.00,8.00) circle (0.43);
  \node[font=\tiny, text=black!55, anchor=west] at (7.43,8.00) {Q5};
  \fill[teal, opacity=0.45] (9.00,9.50) circle (0.27);
  \node[font=\tiny, text=black!55, anchor=west] at (9.27,9.50) {Q6};
  \fill[teal, opacity=0.45] (1.50,3.00) circle (0.41);
  \node[font=\tiny, text=black!55, anchor=west] at (1.91,3.00) {Q7};
  \fill[teal, opacity=0.45] (3.00,1.50) circle (0.31);
  \node[font=\tiny, text=black!55, anchor=west] at (3.31,1.50) {Q8};
  \fill[teal, opacity=0.45] (5.00,3.50) circle (0.55);
  \node[font=\tiny, text=black!55, anchor=west] at (5.55,3.50) {Q9};
  \fill[teal, opacity=0.45] (6.50,2.00) circle (0.35);
  \node[font=\tiny, text=black!55, anchor=west] at (6.85,2.00) {Q10};
  \fill[teal, opacity=0.45] (8.00,4.00) circle (0.44);
  \node[font=\tiny, text=black!55, anchor=west] at (8.44,4.00) {Q11};
  \fill[teal, opacity=0.45] (9.50,1.00) circle (0.25);
  \node[font=\tiny, text=black!55, anchor=west] at (9.75,1.00) {Q12};
  % traiettoria del compromesso efficienza-equità
  \draw[black!45, thick, densely dotted] (5.496,5.029) -- (5.449,5.026) -- (5.401,5.024) -- (5.354,5.021) -- (5.308,5.018) -- (5.261,5.015) -- (5.215,5.013) -- (5.169,5.010) -- (5.124,5.007) -- (5.078,5.005) -- (5.048,5.018) -- (5.032,5.048) -- (5.016,5.077) -- (5.000,5.107) -- (4.983,5.136) -- (4.967,5.165) -- (4.951,5.194) -- (4.935,5.223) -- (4.918,5.253) -- (4.902,5.282) -- (4.886,5.310);
  % vincolo geografico: cerchio della cabina
  \draw[viola, dashed, thick] (7.0,6.0) circle (2.0);
  \node[rectangle, fill=viola, minimum size=2.2mm, inner sep=0] at (7.0,6.0) {};
  \node[star, star points=5, fill=rossomattone, minimum size=3.2mm, inner sep=0pt, label={[font=\scriptsize, text=rossomattone, label distance=2.5mm]190:Weber}] at (4.886,5.310) {};
  \node[star, star points=5, fill=arancio, minimum size=3.2mm, inner sep=0pt, label={[font=\scriptsize, text=arancio, label distance=2.5mm]-55:minimax}] at (5.496,5.029) {};
  \node[star, star points=5, fill=verde, minimum size=3.2mm, inner sep=0pt, label={[font=\scriptsize, text=verde, label distance=2.5mm]120:baricentro}] at (4.897,5.397) {};
  \node[star, star points=5, fill=viola, minimum size=3.2mm, inner sep=0pt, label={[font=\scriptsize, text=viola, label distance=2.5mm]35:Weber vincolato}] at (5.083,5.429) {};
\end{tikzpicture}}
\column{0.5\textwidth}\small
12 quartieri; vincolo: entro 2 km dalla cabina in $(7, 6)$.
\begin{itemize}
  \item Weber $(4{,}89;\, 5{,}31)$: media 3,344 km, max 6,314.
  \item Minimax $(5{,}50;\, 5{,}03)$: media 3,391, max \textbf{5,680}.
  \item Vincolo cabina: $+0{,}2\%$ --- \alert{quasi gratis}.
  \item Frontiera efficienza-equità ($\varepsilon$-constraint) quasi piatta:
        $-0{,}63$ km al più lontano costano 47 m di media.
\end{itemize}
\end{columns}
\end{frame}

% ==================================================================
\section{Ricarica di veicoli elettrici (LP/QP)}

\begin{frame}{Ricarica EV --- il problema}
\textbf{A parole.} Quanta potenza dare a ogni veicolo in ogni ora; obiettivo: minima
spesa; vincoli: energia richiesta entro la partenza, potenza del caricatore quando
collegati, limite del contatore.
\begin{block}{Modello (LP: nessuna variabile binaria!)}
\vspace{-2ex}
\begin{align*}
\min\;& \textstyle\sum_{v \in V}\sum_{t \in T} \pi_t \Delta t\, x_{vt}
\quad\text{s.t.}\;\;
\eta \textstyle\sum_{t \in T} \Delta t\, x_{vt} \ge e_v \;\forall v,\\
& 0 \le x_{vt} \le a_{vt} \bar p_v \;\forall v, t,
\qquad \textstyle\sum_{v \in V} x_{vt} + b_t \le k \;\forall t
\end{align*}
\end{block}
La disponibilità $a_{vt} \in \{0,1\}$ è un \alert{dato}, non una decisione: entra
nel bound. Esempio a mano (2 ore): il duale del fabbisogno $=$ prezzo dell'ora
\alert{marginale}.
\end{frame}

\begin{frame}{Ricarica EV --- costo minimo vs peak shaving}
\begin{columns}
\column{0.58\textwidth}
\begin{tikzpicture}
\begin{axis}[slide, width=\textwidth, height=5.6cm, const plot,
  xlabel={ora}, ylabel={prelievo (kW)}, ymin=0, ymax=130]
\addplot [grigiomedio, dashed] table [col sep=comma, x=ora, y=base] {\figdat{cap10_profili}};
\addplot [teal, very thick] table [col sep=comma, x=ora, y=totcosto] {\figdat{cap10_profili}};
\addplot [arancio, very thick] table [col sep=comma, x=ora, y=totpicco] {\figdat{cap10_profili}};
\addplot [rossomattone, densely dotted, thick] coordinates {(0,120) (23,120)};
\legend{base, costo minimo, peak shaving, limite}
\end{axis}
\end{tikzpicture}
\column{0.42\textwidth}\small
\begin{itemize}
  \item Costo minimo: 20,36~\euro/notte, picco \textbf{103 kW}.
  \item Peak shaving: picco 68 kW; col compromesso costo$+\rho\,$picco lo stesso
        picco costa 22,15~\euro{} (il minimax puro: 27,79!).
  \item Duali $= 0{,}08/\eta$ o $0{,}09/\eta$: l'ora marginale di ogni veicolo.
  \item Capacità minima ammissibile: 68 kW (bisezione).
\end{itemize}
\end{columns}
\begin{block}{Trappola vera}
\texttt{quicksum(x) + base[t] <= C}: il RHS memorizzato è \texttt{C - base[t]} ---
aggiornarlo con \texttt{nuovaC} allenta il vincolo sbagliato.
\end{block}
\end{frame}

% ==================================================================
\section{Code e capacità di servizio (NLP convesso)}

\begin{frame}{Code M/M/1 --- il problema}
\textbf{A parole.} Quanta capacità di servizio comprare; obiettivo: costo di
capacità $+$ valore del tempo dei clienti; vincolo: stabilità $\mu > \lambda$.
\begin{block}{Modello e soluzione chiusa}
\vspace{-1ex}
\[
\min_{\mu > \lambda} \; c\,\mu + h\,\frac{\lambda}{\mu - \lambda}
\qquad\Longrightarrow\qquad
\mu^* = \lambda + \sqrt{h\lambda/c}
\]
``capacità $=$ domanda $+$ \alert{cuscinetto}''; il cuscinetto cresce come
$\sqrt\lambda$ (economie di scala, pooling).
\end{block}
Esempio a mano ($\lambda = 8$, $c = 2$, $h = 4$): $\mu^* = 12$, $\rho = 67\%$,
costo 32~\euro/h; con $\mu = 9$ (``più efficiente'', $\rho = 89\%$): 50~\euro/h.
\end{frame}

\begin{frame}{Code --- il muro dell'utilizzazione e il prezzo di una promessa}
\begin{columns}
\column{0.5\textwidth}
\begin{tikzpicture}
\begin{axis}[slide, width=\textwidth, xlabel={utilizzazione $\rho$ (\%)},
  ylabel={tempo medio $w$ (min)}, ymax=125]
\addplot [teal, very thick] table [col sep=comma, x=rho, y=W_min] {\figdat{cap11_muro}};
\addplot [rossomattone, dash dot] coordinates {(90.2, 0) (90.2, 125)};
\legend{$w(\rho)$, ottimo $90\%$}
\end{axis}
\end{tikzpicture}
\column{0.5\textwidth}\small
Caso di studio: $\lambda = 42$, $c = 3$, $h = 1{,}5$ $\Rightarrow$
$\mu^* = 46{,}6$, $\rho = 90\%$, $w = 13$ min.
\begin{itemize}
  \item Tra $\rho = 90\%$ e $99\%$ l'attesa si moltiplica per \textbf{dieci}:
        conflitto matematico, non organizzativo.
  \item Promessa $w \le w_{\max}$: da 12 a 9 min costa 1,85~\euro/h; da 3 a 2 min
        \alert{28,95}~\euro/h ($dC/dw_{\max} = -c/w_{\max}^2 + h\lambda$).
  \item $\lambda$ incerto in $[36, 48]$: l'ottimo nominale \alert{diverge}
        ($\mu^* < 48$); il robusto ($\mu = 52{,}9$) costa $+13\%$.
\end{itemize}
\end{columns}
\end{frame}

% ==================================================================
\section{Il Newsvendor (LP stocastico)}

\begin{frame}{Newsvendor --- decidere prima di sapere}
\textbf{A parole.} Un'unica quantità $q$, scelta \emph{prima} di osservare la
domanda $D$; obiettivo: minimo costo atteso di eccedenza ($c_o$) e carenza ($c_u$).
\begin{block}{Modello e regola del quantile}
\vspace{-1ex}
\[
\min_{q \ge 0}\; c_o \E[(q - D)^+] + c_u \E[(D - q)^+]
\qquad\Longrightarrow\qquad
F(q^*) = \frac{c_u}{c_u + c_o}
\]
\end{block}
\begin{block}{Esempio a mano ($D \in \{80, 100, 120\}$, $c_u = 9$, $c_o = 4$)}
$C(80) = 180$, $C(100) = 86{,}7$, $C(120) = \textbf{80}$: si ordina il
\alert{massimo}, non la media --- $F(100) = 0{,}67 < \alpha^* = 0{,}692$.
\end{block}
Caso continuo ($\mu = 100$, $\sigma = 20$, $c_u/c_o = 9/4$):
$q^* = F^{-1}(0{,}6923) = 110 \ne 100$.
\end{frame}

\begin{frame}{Newsvendor --- LP a scenari e valore della soluzione stocastica}
\begin{block}{Il trucco delle parti positive (senza binarie)}
\vspace{-2ex}
\[
\min\; c_o \sum_{s \in S} \pi_s o_s + c_u \sum_{s \in S} \pi_s u_s
\quad\text{s.t.}\;\; o_s \ge q - d_s,\;\; u_s \ge d_s - q,\;\; q, o_s, u_s \ge 0
\]
All'ottimo $o_s = (q - d_s)^+$ e $u_s = (d_s - q)^+$: l'obiettivo li schiaccia.
\end{block}
\begin{itemize}
  \item $q$ dell'LP $=$ quantile empirico (108,6 con 600 scenari): è un teorema.
  \item \textbf{VSS} $= 11{,}17$~\euro/ciclo ($-11\%$): il guadagno di modellare
        l'incertezza invece di ordinare ``la media''.
  \item Stabilità: con 30 scenari la $q$ balla di $\pm 3{,}5$ unità tra repliche ---
        gli scenari sono un campione.
  \item Service level: 95\% di non-stockout costa $+45{,}59$~\euro; il fill rate in
        $q^*$ è già il 96\% --- \alert{leggere bene lo SLA}.
\end{itemize}
\end{frame}

\begin{frame}{Newsvendor --- avversione al rischio e multiprodotto}
\begin{columns}
\column{0.52\textwidth}
\begin{tikzpicture}
\begin{axis}[slide, width=\textwidth, xlabel={CVaR$_{0,90}$ del costo (\euro)},
  ylabel={costo medio (\euro)}]
\addplot+ [mark=*, mark size=1.6pt] table [col sep=comma, x=cvar, y=costo_medio]
  {\figdat{cap12_frontiera_cvar}};
\end{axis}
\end{tikzpicture}
\column{0.48\textwidth}\small
\begin{itemize}
  \item Media-CVaR: da $\lambda = 0$ a $1$ si ordina di più (108,6 $\to$ 116,2):
        $+6{,}5$~\euro{} di media comprano $-21{,}4$~\euro{} di coda.
  \item Frontiera ripida all'inizio: molta protezione quasi gratis.
  \item Multiprodotto con budget 1200~\euro{} e domande correlate
        ($\rho = 0{,}7$): duale del budget $-0{,}55$ (rendimento 55\%!);
        con $\rho = 0$ il costo atteso scende del 18\%: la correlazione è un costo.
\end{itemize}
\end{columns}
\end{frame}

% ==================================================================
\section{VaR e CVaR (LP)}

\begin{frame}{VaR e CVaR --- misurare la coda}
$L$ perdita aleatoria, $\alpha \in (0,1)$:
\[
\var_\alpha(L) = \inf\{\eta : \Prob(L \le \eta) \ge \alpha\}
\qquad
\cvar_\alpha(L) = \text{media della coda peggiore di massa } 1 - \alpha
\]
\begin{itemize}
  \item Il CVaR è \alert{convesso} e subadditivo (premia la diversificazione);
        il VaR no.
  \item Esempio a mano (perdite $\{2,4,5,7,12,20\}$, $\alpha = 0{,}8$):
        VaR $= 12$; CVaR $= 12 + \frac{1}{0{,}2}\cdot\frac{1}{6}(20 - 12) =
        \textbf{18,67}$; media $= 8{,}33$: tre numeri, tre storie.
\end{itemize}
\begin{block}{Formulazione lineare di Rockafellar--Uryasev}
\vspace{-1ex}
\[
\min_{\vet x, \eta, \vet\xi}\; \eta + \frac{1}{1 - \alpha}\sum_{s \in S} \pi_s \xi_s
\quad\text{s.t.}\;\; \xi_s \ge \ell_s(\vet x) - \eta,\; \xi_s \ge 0 \;\forall s \in S
\]
All'ottimo $\eta^*$ è un VaR e l'obiettivo è il CVaR: \alert{un solo LP}, entrambe
le misure. ($\eta$ è una variabile \emph{libera}!)
\end{block}
\end{frame}

\begin{frame}{VaR e CVaR --- portafoglio mean-CVaR vs Markowitz}
\begin{columns}
\column{0.52\textwidth}
\begin{tikzpicture}
\begin{axis}[slide, width=\textwidth, ybar, bar width=2.2pt,
  xlabel={perdita mensile (\%)}, ylabel={scenari}, ymin=0]
\addplot [fill=teal!70, draw=none, forget plot] table [col sep=comma, x=centro, y=freq]
  {\figdat{cap13_istogramma}};
\addplot [verde, dashed, thick] coordinates {(-0.67, 0) (-0.67, 30)};
\addplot [arancio, dash dot, thick] coordinates {(3.27, 0) (3.27, 30)};
\addplot [rossomattone, thick] coordinates {(4.89, 0) (4.89, 30)};
\legend{media, VaR$_{90}$, CVaR$_{90}$}
\end{axis}
\end{tikzpicture}
\column{0.48\textwidth}\small
220 scenari con \alert{code grasse} ($t$ di Student), rendimento richiesto 8\%.
\begin{itemize}
  \item Stessa perdita media, ma CVaR$_{90}$: mean-CVaR 4,89\% vs Markowitz 5,31\%
        ($-8\%$).
  \item La varianza penalizza sopra e sotto la media; il CVaR guarda \alert{solo
        dove fa male}.
  \item Con $\alpha = 0{,}99$ e 220 scenari la coda ha 2--3 punti: stime instabili
        --- servono scenari oltre il quantile.
\end{itemize}
\end{columns}
\end{frame}

\begin{frame}{VaR e CVaR --- supply chain a due stadi: il costo della resilienza}
Primo stadio: capacità prenotata $x_a$ dai fornitori (F1 economico ma fragile,
F2 caro ma affidabile). Secondo stadio (per scenario): flussi e shortage.
\begin{center}\small
\begin{tabular}{lcccc}
\toprule
$\lambda$ & prenoto F1 & prenoto F2 & costo medio & CVaR$_{90}$ \\
\midrule
0 (neutrale) & 161,5 & 134,2 & 1.724 & 2.902 \\
0,5          &  92,7 & 212,7 & 1.803 & 2.217 \\
1 (protetto) &  68,5 & 236,9 & 1.906 & 2.139 \\
\bottomrule
\end{tabular}
\end{center}
\begin{block}{Take-away}
Al crescere dell'avversione al rischio la capacità migra verso il fornitore
affidabile: $+79$~\euro{} di costo medio comprano $-685$~\euro{} di CVaR ---
il numero che serve nelle discussioni su dual sourcing e reshoring.
\end{block}
\end{frame}

% ==================================================================
\section{Support Vector Machine (QP)}

\begin{frame}{SVM --- il classificatore come QP}
\textbf{A parole.} Trovare l'iperpiano che separa due classi massimizzando il
margine; le etichette $y_i \in \{-1, +1\}$ sono \alert{dati}, le variabili
($\vet w$, $b$, scarti $\xi_i$) tutte continue.
\begin{block}{Soft margin (primale)}
\vspace{-2ex}
\[
\min_{\vet w, b, \vet\xi}\; \tfrac12 \|\vet w\|^2 + C \sum_{i=1}^n \xi_i
\quad\text{s.t.}\;\;
y_i \Bigl( \sum_{j=1}^{p} w_j x_{ij} + b \Bigr) \ge 1 - \xi_i,\;\; \xi_i \ge 0
\]
margine $= 2/\|\vet w\|$; $C$ = prezzo delle violazioni.
\end{block}
Esempio a mano (2 punti): $\vet w = (0{,}5;\, 0{,}5)$, $b = -1$, margine
$2\sqrt2$ $=$ distanza tra i punti; entrambi support vector.
\end{frame}

\begin{frame}{SVM --- caso di studio: rischio di credito}
90 clienti, due indicatori; al variare di $C$:
\begin{center}\small
\begin{tabular}{lccc}
\toprule
 & $C = 0{,}05$ & $C = 1$ & $C = 20$ \\
\midrule
margine & 2,42 & 1,28 & 0,36 \\
errori di training & 2 & 2 & 0 \\
punti nel margine & 22 & 7 & 3 \\
\bottomrule
\end{tabular}
\end{center}
\begin{itemize}
  \item \textbf{Duale}: $\max \sum_i \alpha_i - \tfrac12 \sum_{i,j} \alpha_i
        \alpha_j y_i y_j \vet x_i\T \vet x_j$, $0 \le \alpha_i \le C$,
        $\sum_i \alpha_i y_i = 0$: stessi $\vet w$, $b$ e valore (dualità forte,
        4,7868).
  \item Complementarietà: $\alpha_i = 0$ (83 punti, irrilevanti);
        $0 < \alpha_i < C$ (2, \alert{sul margine}: da loro si ricava $b$);
        $\alpha_i = C$ (5, dentro il margine).
  \item Classi sbilanciate: costo $10\times$ sugli insolventi $\Rightarrow$
        recall 100\% al prezzo di precision 97\%.
\end{itemize}
\end{frame}

\begin{frame}{SVM --- kernel e SVR}
\begin{columns}
\column{0.5\textwidth}
\begin{tikzpicture}
\begin{axis}[slide, width=\textwidth, height=5.8cm, xlabel={$x_1$}, ylabel={$x_2$},
  xmin=-3.2, xmax=3.2, ymin=-3.2, ymax=3.2]
\addplot [only marks, mark=*, mark size=1.1pt, teal]
  table [col sep=comma, x expr={\thisrow{y}>0 ? \thisrow{x1} : nan}, y=x2] {\figdat{cap14_corona}};
\addplot [only marks, mark=*, mark size=1.1pt, rossomattone]
  table [col sep=comma, x expr={\thisrow{y}<0 ? \thisrow{x1} : nan}, y=x2] {\figdat{cap14_corona}};
\addplot [verde, thick] table [col sep=comma, x=x, y=y] {\figdat{cap14_rbf_g0_7}};
\legend{normale, anomalia, RBF $\gamma = 0{,}7$}
\end{axis}
\end{tikzpicture}
\column{0.5\textwidth}\small
\begin{itemize}
  \item Il duale dipende dai dati solo via $\vet x_i\T \vet x_j$: sostituendolo con
        $K(\vet x, \vet z) = e^{-\gamma\|\vet x - \vet z\|^2}$ si ottengono frontiere
        curve \alert{restando in un QP convesso}.
  \item $\gamma$ piccolo: frontiera regolare; $\gamma = 5$: 61 support vector su 90,
        il modello \alert{memorizza} (overfitting). $C$ e $\gamma$ per
        cross-validation.
  \item \textbf{SVR}: tubo $\pm\varepsilon$ senza penalità; sui dati prezzo-domanda
        recupera $-10{,}86p + 220{,}3$ (vera: $-11p + 220$) usando solo i 14 punti
        fuori dal tubo.
\end{itemize}
\end{columns}
\end{frame}

% ==================================================================
\section{Organizzazione del laboratorio}

\begin{frame}{I quattro laboratori e la consegna}
\begin{center}\small
\begin{tabular}{lll}
\toprule
Lab 1 & Produzione e scorte & LP, duali, verifica dei prezzi ombra \\
Lab 2 & Markowitz & QP, frontiera, fragilità delle stime \\
Lab 3 & Pricing \emph{o} budget & NLP, concavità, KKT numeriche \\
Lab 4 & Progetto a scelta & presentazione manageriale \\
\bottomrule
\end{tabular}
\end{center}
\textbf{Report (max 8 pagine):} problema e ipotesi $\to$ modello (vincoli spiegati)
$\to$ dati $\to$ risultati $\to$ sensitività (protocollo completo) $\to$
\alert{raccomandazione manageriale in dieci righe senza formule}.

\vspace{1ex}
\textbf{Valutazione:} formulazione 30\% $\cdot$ implementazione e verifica 25\%
$\cdot$ sensitività 25\% $\cdot$ interpretazione 20\%.
\end{frame}

\begin{frame}{Gli errori più comuni (checklist)}
\begin{enumerate}
  \item Leggere \texttt{.X} o \texttt{.Pi} senza controllare \texttt{m.Status}.
  \item Dimenticare \texttt{lb=-GRB.INFINITY} sulle variabili libere ($b$, $\eta$).
  \item Usare un prezzo ombra fuori dal suo intervallo di validità.
  \item Sbagliare il segno dei duali nei problemi di minimo.
  \item Aggiornare il RHS di un vincolo con costanti a sinistra.
  \item Ottimizzare un solo obiettivo quando ce ne sono due (minimax puro).
  \item Scegliere gli iperparametri guardando il test set.
  \item Sei cifre decimali da stime che ballano alla seconda.
\end{enumerate}
\end{frame}

\begin{frame}{Materiali}
\begin{itemize}
  \item \textbf{Dispensa completa}: teoria, esempi svolti a mano, casi di studio,
        codice, sensitività, esercizi (con soluzioni per i docenti).
  \item \textbf{Script Python}: uno per capitolo; \texttt{python/esegui\_tutti.py}
        rigenera dati, risultati e figure.
  \item \textbf{Guida al solver}: \texttt{GUIDA\_GUROBI.md}.
  \item \textbf{Dati} in CSV con seed fissi: tutto riproducibile.
\end{itemize}
\vspace{2ex}
\centering
\Large\color{blunotte} Buon laboratorio!
\end{frame}

\end{document}
